\documentclass[11pt]{article}
\usepackage[a4paper,margin=24mm]{geometry}
\usepackage{amsmath,amssymb,amsthm,mathtools}
\usepackage{graphicx}
\usepackage{microtype}
\usepackage[hidelinks]{hyperref}

\newtheorem{theorem}{Theorem}
\newtheorem{proposition}{Proposition}
\newtheorem{corollary}{Corollary}
\newcommand{\C}{\mathbb C}
\newcommand{\A}{\mathcal A}
\newcommand{\supp}{\operatorname{supp}}
\newcommand{\spec}{\sigma}

\title{Spectral invariance on the compact core\\
\large and a nonabelian full subcase of the weighted $L^p$ symmetry conjecture}
\author{Candidate partial-result packet}
\date{August 11, 2026}

\begin{document}
\maketitle

\noindent\textbf{Status.}
Candidate substantial partial result, likely valid pending expert review.

\section{Source conjecture}

Let $G$ be a compactly generated locally compact group of polynomial growth,
$1<p<\infty$, and $1/p+1/q=1$.  A continuous symmetric weight
$\omega\colon G\to[1,\infty)$ satisfies \emph{(LPAlg)} when, up to harmless
positive constants,
\[
 \omega(xy)\leq \omega(x)\omega(y),\qquad
 \omega^{-q}*\omega^{-q}\leq \omega^{-q}.
\]
Then $\A=L^p(G,\omega)$ is a Banach $*$-algebra for convolution and embeds
continuously into $L^1(G)$.  For a relatively compact symmetric generating
neighbourhood $U$, put
\[
 s(n)=\sup_{x\in U^n}\omega(x).
\]
Condition (S) is $s(n)^{1/n}\to1$; on groups of polynomial growth it is
equivalent to the GRS condition \cite{KuznetsovaMolitor,FGL}.

Kuznetsova and Molitor-Braun prove symmetry when $G$ is abelian and, for
general $G$, when $\omega$ is polynomial in Pytlik's sense.  They then ask
whether (LPAlg) plus (S) suffices in general \cite[p.~12]{KuznetsovaMolitor}.

\begin{figure}[ht]
  \centering
  \includegraphics[width=0.92\textwidth]{figures/open_problem_crop.png}
  \caption{The source's Theorem 3.9 and the conjecture immediately following
  it, source PDF page 12.}
\end{figure}

\section{Proof intuition}

Condition (S) is exactly strong enough to remove the weight from high
convolution powers of a compactly supported function.  If $f$ is supported
in $U^m$, then $f^{*n}$ is supported in $U^{mn}$, so its weighted norm costs
at most $s(mn)$.  Young's inequality leaves one $L^p$ factor and measures the
other $n-1$ factors in $L^1$.  Taking $n$th roots kills $s(mn)$ and recovers
the $L^1$ spectral radius.  Applying this to every polynomial in a
self-adjoint $f$ upgrades equality of radii to equality of spectra.

The general completion is delicate because compact truncations need not
preserve spectral radius in a noncommutative Banach algebra.  Two settings
avoid this.  Central elements of a discrete group possess \emph{commuting}
finite-support truncations.  More decisively, if $G=H\times F$ with $H$
abelian and $F$ finite, the whole algebra is, up to an equivalent norm, a
finite direct sum of matrix algebras over the already symmetric abelian
weighted algebra.  This proves the full conjecture for a genuinely
nonabelian family.

\section{The compact core is self-adjoint spectrally invariant}

Write $r_{p,\omega}$ and $r_1$ for spectral radii in $L^p(G,\omega)$ and
$L^1(G)$, using unitizations when needed.

\begin{theorem}[Compact-core spectral invariance]\label{thm:core}
Assume (LPAlg) and condition (S).  If $f\in L^p(G,\omega)$ has compact
support, then
\[
 r_{p,\omega}(f)=r_1(f).
\]
If in addition $f=f^*$, then
\[
 \spec_{L^p(G,\omega)}(f)=\spec_{L^1(G)}(f)\subset\mathbb R.
\]
In particular, the dense $*$-subalgebra $C_c(G)$ is Hermitian with the
spectrum computed in the completed weighted algebra.
\end{theorem}

\begin{proof}
Choose $m$ with $\supp f\subset U^m$.  For $n\geq2$,
\[
 \supp(f^{*n})\subset U^{mn}
\]
and unweighted Young's inequality gives
\begin{align*}
 \|f^{*n}\|_{p,\omega}
 &\leq s(mn)\|f^{*n}\|_p\\
 &\leq s(mn)\|f^{*(n-1)}\|_1\|f\|_p.
\end{align*}
Consequently,
\[
 r_{p,\omega}(f)
 \leq \lim_{n\to\infty}s(mn)^{1/n}
       \|f^{*(n-1)}\|_1^{1/n}\|f\|_p^{1/n}
 =r_1(f),
\]
because $s(mn)^{1/n}=(s(mn)^{1/(mn)})^m\to1$.  Conversely, (LPAlg) and
$\omega^{-q}\in L^1(G)$ give a continuous algebra embedding
$L^p(G,\omega)\hookrightarrow L^1(G)$, hence $r_1(f)\leq r_{p,\omega}(f)$.

Now suppose $f=f^*$ and set $K=\spec_{L^1(G)}(f)$.  The group $G$ has
polynomial growth, so $L^1(G)$ is symmetric by Losert's theorem; hence
$K\subset\mathbb R$.  For every polynomial $P$ with $P(0)=0$, the element
$P(f)$ remains compactly supported.  The first part and spectral mapping give
\[
 r_{p,\omega}(P(f))=r_1(P(f))=\max_{z\in K}|P(z)|.
\]
If $\lambda\in\spec_{L^p(G,\omega)}(f)$, spectral mapping therefore yields
\[
 |P(\lambda)|\leq\max_{z\in K}|P(z)|
 \quad\text{for all such }P.
\]
A compact subset of the real line is polynomially convex.  In the nonunital
case $0\in K$, so polynomials vanishing at zero still separate every point
outside $K$.  Thus $\lambda\in K$.  The reverse spectral inclusion follows
from the algebra embedding, proving equality.
\end{proof}

The theorem is stronger than a mere density statement: spectra are computed
in the completed weighted algebra.  It also isolates the remaining issue:
spectral radius is not generally continuous under noncommuting norm
approximation.

\section{A completion upgrade for discrete central elements}

\begin{proposition}\label{prop:center}
Suppose in addition that $G$ is discrete.  If $z$ is central in
$L^p(G,\omega)$, then
\[
 r_{p,\omega}(z)=r_1(z).
\]
If $z=z^*$, then its spectra in $L^p(G,\omega)$ and $L^1(G)$ coincide and
are real.
\end{proposition}

\begin{proof}
A central function on a discrete group is constant on conjugacy classes.
Since $\omega\geq1$, a weighted $\ell^p$ function must vanish on every
infinite conjugacy class.  It follows that $z$ has finite-support central
truncations $z_N$ with $\|z-z_N\|_{p,\omega}\to0$.  Put $h_N=z-z_N$.
All three elements commute.  Spectral radius is subadditive on commuting
elements, Theorem~\ref{thm:core} applies to $z_N$, and the embedding into
$\ell^1(G)$ has some norm $C$.  Hence
\begin{align*}
 r_{p,\omega}(z)
 &\leq r_{p,\omega}(z_N)+\|h_N\|_{p,\omega}\\
 &=r_1(z_N)+\|h_N\|_{p,\omega}\\
 &\leq r_1(z)+\|h_N\|_1+\|h_N\|_{p,\omega}\\
 &\leq r_1(z)+(C+1)\|h_N\|_{p,\omega}.
\end{align*}
Letting $N\to\infty$ and using the reverse inequality proves equality.
The same radius equality applies to every polynomial in a central $z$.
When $z=z^*$, polynomial separation of its real $L^1$ spectrum, as in
Theorem~\ref{thm:core}, gives equality of spectra.
\end{proof}

\section{A full nonabelian subcase}

\begin{theorem}[Finite direct factors]\label{thm:finite}
Let $H$ be a compactly generated locally compact abelian group of polynomial
growth, let $F$ be a finite group, and put $G=H\times F$.  If
$(G,\omega)$ satisfies (LPAlg) and condition (S), then
\[
 L^p(G,\omega)
\]
is a symmetric Banach $*$-algebra.  In particular, this proves the source
conjecture for nonabelian groups such as $\mathbb R\times S_3$, for every
admissible GRS weight and not only for Pytlik-polynomial weights.
\end{theorem}

\begin{proof}
Normalize Haar measure on the finite factor arbitrarily and define
\[
 \omega_H(h)=\omega(h,e_F),\qquad B=L^p(H,\omega_H).
\]
For $t\in F$, submultiplicativity gives
\[
 \frac{\omega_H(h)}{\omega(e_H,t^{-1})}
 \leq \omega(h,t)
 \leq \omega_H(h)\omega(e_H,t).
\]
Because $F$ is finite, these constants are uniform in $t$.  Thus, as a
Banach space with equivalent norms,
\[
 L^p(H\times F,\omega)=\bigoplus_{t\in F} B.
\]

The restriction $\omega_H$ is symmetric and submultiplicative.  In the
convolution inequality for $\omega^{-q}$ on $H\times F$, retain only the
$t=e_F$ summand.  This gives
\[
 \omega_H^{-q}*_H\omega_H^{-q}\leq C\omega_H^{-q}
\]
for a constant depending only on the Haar normalization of $F$.  After the
standard scalar renormalization, $(H,\omega_H)$ satisfies (LPAlg).
Condition (S) also restricts from $H\times F$ to $H$: use a generating set
$V\times F$, where $V$ generates $H$.  The abelian theorem of
\cite{KuznetsovaMolitor} therefore shows that $B$ is symmetric.

Convolution on the direct product identifies the Banach $*$-algebra, up to
the equivalent norm already noted, with the finite tensor product
\[
 B\otimes\C[F].
\]
Finite-group Fourier theory gives a $*$-algebra decomposition
\[
 \C[F]\cong\bigoplus_{\pi\in\widehat F}M_{d_\pi}(\C),
\]
and hence
\[
 B\otimes\C[F]
 \cong\bigoplus_{\pi\in\widehat F}M_{d_\pi}(B).
\]
Symmetry of Banach $*$-algebras is preserved by finite matrix amplification
and finite direct sums (the Leptin--Wichmann matrix theorem; see
\cite{Palmer}).  Therefore $L^p(H\times F,\omega)$ is symmetric.
\end{proof}

\section{Why the general completion remains open}

Hulanicki's transfer lemma requires equality of spectral radii for
\emph{every} element of $\A$ \cite[Lemma~3.1]{FGL}; an arbitrary dense core
does not suffice, and spectral radius need not behave continuously under
noncommuting compact truncations.  The weighted-$L^1$ annular proof also
requires a weighted first moment, whereas $L^p(G,\omega)$ supplies only a
weighted $p$-moment and an unweighted first moment.  Eight focused attempts
confirmed this completion barrier.  The general result is not claimed; a
further upgrade needs a new all-elements inverse-closedness or representation
theorem.

\section{Verification and novelty status}

No numerical evidence is used.  The support estimate, Hulanicki hypothesis,
finite-factor coordinates, and Haar-normalization constants were audited
directly.  On August 11, 2026, bounded index and web searches by exact title,
quoted question, authors, weighted-$L^p$, GRS, and condition (S) found no
later general solution or statement of Theorem~\ref{thm:finite}.  Novelty
confidence is moderate pending specialist review, especially of the
nonunital polynomial separation and finite matrix reduction.

{\footnotesize
\begin{thebibliography}{9}
\setlength{\itemsep}{1pt}

\bibitem{KuznetsovaMolitor}
Y. Kuznetsova and C. Molitor-Braun,
\emph{Harmonic analysis of weighted $L^p$-algebras},
Expositiones Mathematicae \textbf{30} (2012), 124--153;
arXiv:1103.3610.

\bibitem{FGL}
G. Fendler, K. Gr\"ochenig, and M. Leinert,
\emph{Symmetry of weighted $L^1$-algebras and the GRS-condition},
Bull. London Math. Soc. \textbf{38} (2006), 625--635.

\bibitem{Palmer}
T. W. Palmer,
\emph{Banach Algebras and the General Theory of $*$-Algebras, Vol. II},
Cambridge University Press, 2001.

\end{thebibliography}
}

\end{document}
